% ============================================================
%  Chapter 2 – State of the Art
%  PsychPlatform – Academic Report
%  Generated from state_of_the_art.md.resolved
% ============================================================

\chapter{State of the Art}
\label{chap:soa}

% ─────────────────────────────────────────────────────────────
\section*{Introduction}
\addcontentsline{toc}{section}{Introduction}
% ─────────────────────────────────────────────────────────────

The preceding chapter established the general context and objectives of this
project: the design and development of an AI-assisted psychological consultation
platform, hereafter referred to as \emph{PsychPlatform}. This platform aims to
connect patients in need of psychological support with verified, licensed
psychologists through a structured and intelligent digital environment. The
system integrates role-based access control, an AI-powered conversational agent,
automated clinical summarisation, real-time communication, geospatial
practitioner discovery, and a document-based credential verification workflow.

The present chapter constitutes a thorough review of the scientific domains
underpinning the project and of the existing solutions that operate within
comparable problem spaces. The purpose of this review is twofold. First, it
provides the conceptual and theoretical foundations necessary to understand the
choices made during the design and development of the system. Second, it situates
the proposed platform within a landscape of existing tools and services,
identifying the gaps and limitations that motivate and justify the specific
contributions of this work.

The chapter is organised as follows. Section~\ref{sec:domains} surveys the
scientific domains and key concepts relevant to the project.
Section~\ref{sec:existing} presents an analysis of existing platforms and
solutions. Section~\ref{sec:comparison} provides a structured comparative
analysis. Section~\ref{sec:limitations} identifies the limitations of current
systems. Section~\ref{sec:opportunities} articulates the opportunities that
remain unaddressed and explains how the proposed platform responds to them.

% ─────────────────────────────────────────────────────────────
\section{Scientific Domains and Key Concepts}
\label{sec:domains}
% ─────────────────────────────────────────────────────────────

The development of PsychPlatform draws upon a broad and interdisciplinary set of
scientific fields. This section introduces each major domain and explains the
specific concepts within that domain that are directly relevant to the
architecture and functionality of the system.

% -----------------------------------------------------------
\subsection{Web Engineering and Full-Stack Development}
\label{subsec:web-eng}
% -----------------------------------------------------------

Web engineering is the discipline concerned with the systematic application of
engineering principles to the design, development, testing, and maintenance of
web-based software systems \cite{Murugesan2001}. Modern web engineering has
evolved significantly from the era of static HTML documents to the present
paradigm of dynamic, data-driven, reactive applications served over distributed
architectures.

\subsubsection{The Client-Server Architecture}

The client-server model is the foundational architectural paradigm of modern web
applications. Under this model, a \emph{client} (typically a browser running a
JavaScript application) communicates with a \emph{server} (a process running on
remote hardware) over the HTTP or HTTPS protocol. The server is responsible for
processing business logic, authenticating requests, and managing persistent data.
The client is responsible for rendering the user interface and managing localised
application state.

PsychPlatform follows this paradigm strictly, with the client implemented in
React and the server implemented in Node.js with Express.

\begin{figure}[H]
  \centering
  % INSERT FIGURE: Diagram illustrating the client-server communication model
  % in PsychPlatform. The figure should show a browser (React SPA) on the left
  % sending HTTP/HTTPS requests to a Node.js/Express server on the right, with
  % arrows indicating the request-response cycle. MongoDB should appear behind
  % the server as the persistence layer. JWT tokens should be shown on the
  % request arrow to illustrate stateless authentication.
  \rule{0.85\textwidth}{0.4pt}\\[6pt]
  {\small\itshape [Figure placeholder — Client-Server Architecture of PsychPlatform]}
  \rule{0.85\textwidth}{0.4pt}
  \caption{High-level client-server communication model in PsychPlatform.
           The React single-page application communicates with the Node.js/Express
           backend exclusively over HTTPS, attaching a JWT token on each protected
           request. MongoDB serves as the persistence layer on the server side.}
  \label{fig:client-server}
\end{figure}

\subsubsection{The MERN Stack}

The MERN stack is a collection of four open-source technologies — MongoDB,
Express.js, React, and Node.js — that together form a comprehensive
JavaScript-based ecosystem for building full-stack web applications. Its
principal advantage is the uniformity of the programming language (JavaScript
and its superset TypeScript) across both the client and server tiers, which
reduces context-switching for developers and facilitates the sharing of
validation logic, data models, and utility functions.

\begin{itemize}[leftmargin=1.5em, itemsep=4pt]

  \item \textbf{MongoDB} is a document-oriented NoSQL database that stores data
        as flexible JSON-like documents (BSON). Its schemaless nature allows for
        rapid, iterative development of data models, which is particularly
        valuable in early-stage projects where requirements evolve frequently.
        In PsychPlatform, MongoDB stores all persistent entities: users,
        psychologist profiles, sessions, messages, chatbot histories, ratings,
        calendar slots, and clinical summaries.

  \item \textbf{Express.js} is a minimal and unopinionated web application
        framework for Node.js. It provides routing, middleware composition, and
        HTTP utility methods, enabling the development of structured RESTful
        APIs. PsychPlatform's backend exposes a comprehensive REST API organised
        into domain-specific route groups (authentication, sessions, chatbot,
        psychologist management, admin, notifications, etc.).

  \item \textbf{React} is a declarative, component-based JavaScript library
        developed by Meta for building user interfaces \cite{Fedosejev2015}.
        React's virtual DOM diffing algorithm allows for efficient, incremental
        updates to the user interface without full page reloads. PsychPlatform's
        frontend is a single-page application (SPA) built with React, using
        React Router for client-side navigation.

  \item \textbf{Node.js} is a JavaScript runtime built on the V8 engine that
        allows JavaScript to be executed outside of a browser, enabling
        server-side development. Its event-driven, non-blocking I/O model makes
        it particularly well-suited to I/O-intensive applications that handle
        many concurrent connections, as is the case for a consultation platform
        managing real-time sessions.

\end{itemize}

\begin{figure}[H]
  \centering
  % INSERT FIGURE: A four-quadrant or layered diagram of the MERN stack.
  % Each quadrant should contain the technology logo placeholder and a brief
  % label: MongoDB (Database), Express.js (API Layer), React (UI Layer),
  % Node.js (Runtime). Arrows should connect layers to show data flow from
  % the database up through the API to the React frontend.
  \rule{0.85\textwidth}{0.4pt}\\[6pt]
  {\small\itshape [Figure placeholder — MERN Stack Layer Diagram]}
  \rule{0.85\textwidth}{0.4pt}
  \caption{The MERN stack as used in PsychPlatform. Each technology occupies a
           distinct architectural tier: MongoDB at the data layer, Express.js at
           the API layer, Node.js as the server runtime, and React as the
           client-side presentation layer. JavaScript is the shared language
           across all tiers.}
  \label{fig:mern-stack}
\end{figure}

\subsubsection{RESTful API Design}

Representational State Transfer (REST) is an architectural style for distributed
hypermedia systems \cite{Fielding2000}. A RESTful API is one that adheres to REST
constraints, including stateless communication, a uniform interface (typically
CRUD operations mapped to HTTP verbs: \texttt{GET}, \texttt{POST}, \texttt{PUT},
\texttt{DELETE}), and resource-based URIs. RESTful APIs have become the de facto
standard for web service design due to their simplicity, scalability, and
compatibility with a wide range of clients.

PsychPlatform exposes a RESTful API with endpoints organised by resource
(e.g., \texttt{/api/sessions}, \texttt{/api/chatbot},
\texttt{/api/psychologists}, \texttt{/api/ratings}), each subject to role-based
access control enforced through JWT middleware.

\subsubsection{Authentication and Authorisation}

Authentication is the process of verifying the identity of a user or system.
Authorisation is the process of determining what an authenticated entity is
permitted to do. These two mechanisms are distinct but complementary.

JSON Web Tokens (JWT) are an open standard (RFC 7519) for securely transmitting
claims between parties as a digitally signed JSON object. JWTs are stateless:
the server does not maintain a session table; instead, the token itself carries
the user's identity and role, and its integrity is guaranteed by a cryptographic
signature. Upon a valid login, PsychPlatform issues a JWT encoding the user's
identifier and role (\texttt{patient}, \texttt{psychologist}, or
\texttt{admin}). This token is attached to all subsequent API requests in the
\texttt{Authorization} header and verified by the server's authentication
middleware before any protected endpoint is accessed.

Password security is enforced through \texttt{bcrypt} hashing. Bcrypt is a
computationally expensive adaptive hashing function designed specifically for
password storage, incorporating a work factor to resist brute-force attacks
\cite{Provos1999}.

\subsubsection{Real-Time Communication with WebSockets}

The HTTP request-response protocol is fundamentally unidirectional: the client
initiates a request, and the server responds. This model is insufficient for
real-time applications where the server must push data to connected clients
asynchronously, such as in live messaging systems.

WebSocket is a protocol (RFC~6455) that establishes a persistent, full-duplex
communication channel over a single TCP connection. Socket.IO is a JavaScript
library that abstracts over WebSocket (with fallbacks to HTTP long-polling when
WebSocket is unavailable), providing room-based broadcasting, event-driven APIs,
and automatic reconnection logic. PsychPlatform uses Socket.IO to power its
real-time session messaging feature, allowing psychologists and patients to
exchange messages instantaneously within a shared room identified by the session
identifier.

\begin{figure}[H]
  \centering
  % INSERT FIGURE: Side-by-side comparison of HTTP polling vs. WebSocket.
  % Left side: classic request-response cycle with multiple separate HTTP
  % requests shown as discrete arrows. Right side: a single persistent
  % bi-directional WebSocket connection with continuous two-way arrows.
  % A Socket.IO room (labelled with a session ID) should be shown in the
  % centre linking a patient client and a psychologist client.
  \rule{0.85\textwidth}{0.4pt}\\[6pt]
  {\small\itshape [Figure placeholder — HTTP vs.\ WebSocket Communication Model]}
  \rule{0.85\textwidth}{0.4pt}
  \caption{Comparison of HTTP request-response polling (left) and the persistent
           WebSocket/Socket.IO channel (right) used in PsychPlatform for real-time
           session messaging. Both the patient and the psychologist client join
           a Socket.IO room keyed by the session identifier.}
  \label{fig:websocket}
\end{figure}

\subsubsection{Internationalisation (i18n)}

Internationalisation refers to the engineering practice of designing a software
system to support multiple languages and locales without requiring changes to the
core codebase \cite{W3C2005}. The \texttt{react-i18next} library, built on top
of the i18next framework, provides a declarative API for declaring translatable
strings using lookup keys and for switching languages at runtime. PsychPlatform
supports three languages — English, French, and Arabic — including right-to-left
(RTL) layout rendering for Arabic, which requires CSS-level direction reversals
for all directionally sensitive UI components.

% -----------------------------------------------------------
\subsection{Artificial Intelligence and Machine Learning}
\label{subsec:ai-ml}
% -----------------------------------------------------------

Artificial Intelligence (AI) encompasses a broad field of computer science
concerned with the development of systems that exhibit behaviours associated with
human intelligence, such as reasoning, learning, natural language understanding,
and perception \cite{Russell2020}. Machine Learning (ML) is a subfield of AI in
which systems learn patterns from data rather than being programmed with explicit
rules \cite{Mitchell1997}.

Within PsychPlatform, AI is employed not as a peripheral feature but as a core
architectural component. The platform integrates AI at three distinct functional
layers:
\begin{enumerate}[label=(\arabic*), itemsep=2pt]
  \item the patient-facing conversational agent;
  \item the automated clinical summarisation and emotional analysis engine;
  \item the credential verification assistant for psychologist onboarding.
\end{enumerate}

\subsubsection{The Role of AI in Mental Health Applications}

The application of AI to mental health support is a rapidly growing area of
research and commercial development. Broadly speaking, AI in this domain serves
two non-exclusive roles. The first is \emph{screening and triage}: AI systems
analyse patient-provided data (text, voice, physiological signals) to detect
indicators of mental health conditions and assess severity. The second is
\emph{therapeutic augmentation}: AI systems supplement, extend, or prepare the
patient for human-led therapeutic sessions by providing a safe, non-judgmental
conversational space.

PsychPlatform belongs to the second category. Its chatbot is explicitly designed
not to diagnose or provide medical advice, but to serve as a preparatory and
expressive tool that lowers the barrier to formal psychological consultation.

% -----------------------------------------------------------
\subsection{Natural Language Processing}
\label{subsec:nlp}
% -----------------------------------------------------------

Natural Language Processing (NLP) is the branch of AI concerned with enabling
computers to understand, interpret, generate, and respond to human language
\cite{JurafskyMartin2023}. NLP spans a wide spectrum of tasks, from low-level
morphological analysis (tokenisation, stemming, lemmatisation) to high-level
semantic understanding (question answering, sentiment analysis, discourse
modelling).

Key NLP tasks relevant to PsychPlatform include:

\begin{itemize}[leftmargin=1.5em, itemsep=6pt]

  \item \textbf{Sentiment Analysis.} The computational identification of
        subjective information — opinions, emotions, and attitudes — from
        textual data. In the context of this platform, sentiment analysis is
        used to assess the emotional tone of a patient's conversational messages
        and to infer a sentiment trend (improving, stable, or declining) over
        time.

  \item \textbf{Emotion Detection.} A more fine-grained form of sentiment
        analysis that classifies text according to discrete emotional categories
        (e.g., anxiety, sadness, anger, joy). PsychPlatform computes emotional
        scores across four dimensions — anxiety, sadness, anger, and positivity
        — for each patient's session.

  \item \textbf{Text Summarisation.} The automatic production of a condensed
        version of a longer text that preserves the essential information.
        Abstractive summarisation, as opposed to extractive summarisation,
        generates novel sentences rather than selecting and concatenating
        existing ones. PsychPlatform employs abstractive summarisation to
        produce clinical summaries (\texttt{rawSummary}) from the full
        transcript of a patient's chatbot conversation.

  \item \textbf{Key Theme Extraction.} The identification of the most salient
        topics or subjects discussed in a body of text (\texttt{keyThemes}).
        This assists the psychologist in rapidly identifying the principal
        concerns a patient has raised before a scheduled session.

  \item \textbf{Multilingual Processing.} NLP at scale must handle text in
        multiple languages. PsychPlatform's chatbot is designed to detect the
        language of the patient's input automatically and respond accordingly,
        supporting Arabic, French, and English within a single conversational
        pipeline.

\end{itemize}

% -----------------------------------------------------------
\subsection{Large Language Models}
\label{subsec:llm}
% -----------------------------------------------------------

Large Language Models (LLMs) are a class of neural network-based AI systems
trained on massive corpora of text with the objective of modelling the
probability distribution of token sequences. The transformer architecture
\cite{Vaswani2017}, with its self-attention mechanism, is the foundational
building block of all modern LLMs. The pretraining phase exposes the model to
billions of text tokens, enabling it to acquire broad linguistic and world
knowledge. Models such as GPT-4 \cite{OpenAI2023}, LLaMA~3 \cite{MetaAI2024},
Claude \cite{Anthropic2024}, and Gemini \cite{GoogleDeepMind2024} have
demonstrated remarkable capabilities in instruction following, multi-turn
dialogue, code generation, and complex reasoning.

\begin{figure}[H]
  \centering
  % INSERT FIGURE: Simplified diagram of the Transformer self-attention mechanism.
  % Show an input token sequence at the bottom, multi-head self-attention blocks
  % in the middle (with Query/Key/Value arrows), and an output probability
  % distribution at the top. Annotate the pre-training (large unlabelled corpus)
  % phase on the left and the fine-tuning / instruction-following phase on the
  % right. Highlight the two LLaMA 3 variants used in PsychPlatform (70B and 8B)
  % in a callout box.
  \rule{0.85\textwidth}{0.4pt}\\[6pt]
  {\small\itshape [Figure placeholder — Transformer Architecture and LLM Pre-training Overview]}
  \rule{0.85\textwidth}{0.4pt}
  \caption{Simplified overview of the transformer-based LLM architecture
           (\citealt{Vaswani2017}). Input tokens are projected into Query, Key,
           and Value matrices; self-attention scores determine contextual
           representations; and a softmax output head produces the next-token
           probability distribution. PsychPlatform employs two LLaMA~3 variants
           served via the Groq inference API.}
  \label{fig:transformer}
\end{figure}

PsychPlatform integrates two LLaMA~3 model variants served through the Groq
inference API:

\begin{itemize}[leftmargin=1.5em, itemsep=6pt]

  \item \textbf{LLaMA~3.3 70B Versatile.} A large, general-purpose model used
        for the patient-facing conversational chatbot and for the platform
        navigation assistant. Its large parameter count enables nuanced,
        contextually appropriate, and empathetic dialogue generation.

  \item \textbf{LLaMA~3.1 8B Instant.} A smaller, faster model used for
        structured JSON generation tasks such as clinical summarisation and
        credential analysis. Its lower latency is advantageous for tasks where
        response speed and output format compliance are more critical than
        response depth.

\end{itemize}

The use of LLMs in PsychPlatform represents a significant departure from earlier
chatbot paradigms based on rule-based systems or retrieval-based models.
Rule-based systems respond according to predefined scripts and decision trees;
they are robust but brittle, failing when user input deviates from anticipated
patterns. Retrieval-based models select responses from a fixed candidate pool,
offering slightly greater flexibility but still constrained by the quality and
scope of the response database. LLMs, by contrast, generate responses
dynamically, enabling them to handle the full unpredictability of natural human
dialogue \cite{Brown2020}.

% -----------------------------------------------------------
\subsection{Prompt Engineering}
\label{subsec:prompt}
% -----------------------------------------------------------

\begin{figure}[H]
  \centering
  % INSERT FIGURE: Pipeline diagram of the LLM prompt engineering workflow
  % used in PsychPlatform. The pipeline should read left-to-right:
  % (1) User message input  →  (2) System prompt construction (role + guardrails
  % + format constraints + temperature setting)  →  (3) Groq API call  →
  % (4) Raw LLM output  →  (5) parseJsonFromModel() cleaning & parsing  →
  % (6) Structured JSON output (dominantEmotion, urgencyScore, etc.).
  % Use colour-coded boxes for each stage.
  \rule{0.85\textwidth}{0.4pt}\\[6pt]
  {\small\itshape [Figure placeholder — Prompt Engineering Pipeline in PsychPlatform]}
  \rule{0.85\textwidth}{0.4pt}
  \caption{End-to-end prompt engineering pipeline in PsychPlatform. A structured
           system prompt — comprising role assignment, behavioural guardrails,
           output format constraints, and a calibrated temperature value — wraps
           the user message before submission to the Groq inference API. The raw
           model response is then cleaned and parsed by \texttt{parseJsonFromModel}
           to yield a reliable, structured JSON object.}
  \label{fig:prompt-pipeline}
\end{figure}

Prompt engineering is the discipline of designing, structuring, and iterating
upon the textual instructions (prompts) provided to an LLM in order to elicit
specific, reliable, and appropriately constrained outputs \cite{White2023}.
Since LLMs are highly sensitive to the phrasing and structure of their input,
prompt engineering is a critical engineering skill in any application that relies
on LLM inference.

PsychPlatform employs several prompt engineering strategies:

\begin{itemize}[leftmargin=1.5em, itemsep=6pt]

  \item \textbf{System Prompt Role Assignment.} Each LLM call is prefixed with
        a system prompt that assigns a specific role to the model (e.g.,
        ``\textit{You are a compassionate listening assistant}'', ``\textit{You
        are a clinical assistant}'', ``\textit{You are an assistant that
        verifies psychologist credentials}''). This shapes the model's output
        register, vocabulary, and behavioural boundaries.

  \item \textbf{Output Format Constraints.} For structured tasks such as
        clinical summarisation, the system prompt explicitly instructs the model
        to return only a valid JSON object with precisely defined fields. This
        is necessary because LLMs by default produce free-form text, and
        consuming applications require machine-readable, predictable output.

  \item \textbf{Behavioural Guardrails.} The chatbot's system prompt includes
        explicit prohibitions (``Never provide medical diagnoses or suggest
        medications'', ``Never rush the patient'') and instructions for edge
        cases (``Always detect the language the patient is writing in and
        respond in the SAME language''), which serve as safety and quality
        constraints.

  \item \textbf{Temperature Control.} The \texttt{temperature} parameter
        governs the randomness of the model's output distribution. A
        temperature of \texttt{0.3} is used for analytical and structured tasks
        (summarisation, verification) to ensure consistent and conservative
        outputs, while \texttt{0.7} is used for conversational generation to
        introduce natural variation and prevent repetitive phrasing.

  \item \textbf{Robust Output Parsing.} In practice, LLMs may deviate from the
        requested output format by wrapping JSON in Markdown code fences or
        inserting explanatory text. PsychPlatform implements a parsing function
        (\texttt{parseJsonFromModel}) that strips known formatting artefacts
        before applying \texttt{JSON.parse}, and falls back to
        boundary-scanning the text for the first and last brace characters when
        standard parsing fails.

\end{itemize}

% -----------------------------------------------------------
\subsection{Sentiment Analysis and Emotional State Modelling}
\label{subsec:sentiment}
% -----------------------------------------------------------

Sentiment analysis occupies a central functional role in PsychPlatform through
the generation of clinical summaries. When a patient's chatbot session is
concluded, the full conversation transcript is submitted to the LLM with a
prompt requesting a structured assessment of the patient's emotional state. The
output includes the following structured fields:

\begin{itemize}[leftmargin=1.5em, itemsep=4pt]
  \item \textbf{Dominant Emotion} (\texttt{dominantEmotion}): A single-word
        classification of the most prevalent emotional state expressed in the
        conversation.
  \item \textbf{Urgency Score} (\texttt{urgencyScore}): A numerical score from
        1 to 5 representing the clinical urgency of the patient's situation,
        where higher values indicate greater need for immediate professional
        attention.
  \item \textbf{Sentiment Trend} (\texttt{sentimentTrend}): A categorical
        assessment of whether the patient's emotional state has been improving,
        stable, or declining over the course of the conversation.
  \item \textbf{Key Themes} (\texttt{keyThemes}): An array of the main topics
        or concerns raised by the patient.
  \item \textbf{Raw Summary} (\texttt{rawSummary}): A free-text narrative
        summary of the full conversation.
\end{itemize}

This information is subsequently made available to the assigned psychologist via
the patient detail view, enabling evidence-based, personalised consultation
preparation.

Furthermore, the platform maintains an \texttt{EmotionalIndicator} model that
stores per-session quantitative scores across four emotional dimensions (anxiety,
sadness, anger, positivity) on a normalised scale of 0 to 100. These indicators
allow for longitudinal tracking of a patient's emotional trajectory across
multiple sessions, supporting data-driven insight for the supervising
psychologist.

\begin{figure}[H]
  \centering
  % INSERT FIGURE: Two-part figure.
  % Left panel: a sample JSON block showing the clinical summary output
  % structure (dominantEmotion, urgencyScore 1-5, sentimentTrend, keyThemes[],
  % rawSummary) as it is returned by the LLM and stored in the database.
  % Right panel: a line chart (multi-series) showing a hypothetical patient's
  % EmotionalIndicator scores (anxiety, sadness, anger, positivity) on the
  % Y-axis (0-100) plotted across five consecutive sessions on the X-axis.
  \rule{0.85\textwidth}{0.4pt}\\[6pt]
  {\small\itshape [Figure placeholder — Clinical Summary Output and Longitudinal Emotional Tracking]}
  \rule{0.85\textwidth}{0.4pt}
  \caption{Left: example structured clinical summary JSON produced by the LLM
           after a patient's chatbot session, highlighting the five output fields
           stored in the database. Right: illustrative longitudinal
           \texttt{EmotionalIndicator} chart for a patient across five
           consecutive sessions, showing per-session scores for anxiety, sadness,
           anger, and positivity on a normalised 0--100 scale.}
  \label{fig:emotion-tracking}
\end{figure}

% -----------------------------------------------------------
\subsection{Geospatial Computing and Location-Based Services}
\label{subsec:geo}
% -----------------------------------------------------------

Geospatial computing involves the processing of geographic information to solve
location-based problems. In the context of web applications, location-based
services (LBS) allow users to discover entities of interest within a defined
geographic proximity.

PsychPlatform integrates geospatial search to enable patients to locate
psychologists near their physical location. The \texttt{Psychologist} schema
stores a GeoJSON \texttt{Point} object containing latitude and longitude
coordinates, and a \texttt{2dsphere} index is created on this field in MongoDB.
This index enables MongoDB's geospatial query operators (e.g., \texttt{\$near},
\texttt{\$geoWithin}) to efficiently compute proximity-based search results,
sorting psychologists by their distance from the patient's coordinates. This
functionality is particularly relevant in contexts where patients prefer or
require in-person consultations when available.

\begin{figure}[H]
  \centering
  % INSERT FIGURE: Map-style diagram illustrating geospatial proximity search.
  % Show a stylised city map with a patient marker at the centre and several
  % psychologist markers at varying distances. Draw concentric radius rings
  % (e.g., 5 km, 10 km, 20 km) around the patient. Psychologists within each
  % ring should be numbered in order of proximity. A MongoDB 2dsphere query
  % snippet ($near operator) should appear as an annotation box.
  \rule{0.85\textwidth}{0.4pt}\\[6pt]
  {\small\itshape [Figure placeholder — Geospatial Proximity Search using MongoDB 2dsphere Index]}
  \rule{0.85\textwidth}{0.4pt}
  \caption{Geospatial proximity search in PsychPlatform. A patient's GPS
           coordinates are used as the origin of a \texttt{\$near} query against
           the \texttt{2dsphere}-indexed \texttt{location} field of the
           \texttt{Psychologist} collection. Results are returned in ascending
           order of distance, enabling the patient to discover the nearest
           available practitioners.}
  \label{fig:geospatial}
\end{figure}

% -----------------------------------------------------------
\subsection{Document Parsing and AI-Assisted Credential Verification}
\label{subsec:cred}
% -----------------------------------------------------------

The onboarding of psychologists to the platform requires validation of
professional credentials to protect patients from unqualified practitioners.
PsychPlatform implements an automated preliminary screening workflow that uses
AI to assist the administrator in this process.

When a psychologist submits their CV and diploma as PDF documents, the server
extracts the textual content of each document using the \texttt{pdf-parse}
library. The extracted text is then submitted to an LLM with a structured prompt
requesting an assessment of the candidate's qualifications, estimated years of
experience, identified specialisations, the apparent legitimacy of the diploma,
and an overall recommendation (Approve, Review, or Reject). The resulting
AI-generated summary (\texttt{aiVerificationSummary}) is stored in the
psychologist's profile and displayed to the administrator in the pending
verifications panel, alongside the document files themselves. The final decision
to approve or reject the psychologist is made by the human administrator,
preserving human oversight in this safety-critical process.

This workflow reflects a broader trend in AI-assisted decision-making: using AI
to augment and accelerate human judgment rather than to fully automate
consequential decisions \cite{Dellermann2019}.

\begin{figure}[H]
  \centering
  % INSERT FIGURE: Sequential flowchart of the credential verification pipeline.
  % Steps (left to right or top to bottom):
  % (1) Psychologist uploads CV + Diploma PDF
  %     ↓
  % (2) Server: pdf-parse extracts raw text
  %     ↓
  % (3) LLM call with structured credential-analysis prompt
  %     ↓
  % (4) AI generates aiVerificationSummary JSON
  %     { qualifications, experience, specialisations, recommendation }
  %     ↓
  % (5) Summary stored in Psychologist profile
  %     ↓
  % (6) Admin reviews: [Approve] / [Reject] (human-in-the-loop)
  % Use colour-coded decision diamond for step 6.
  \rule{0.85\textwidth}{0.4pt}\\[6pt]
  {\small\itshape [Figure placeholder — AI-Assisted Credential Verification Workflow]}
  \rule{0.85\textwidth}{0.4pt}
  \caption{Step-by-step AI-assisted credential verification workflow in
           PsychPlatform. PDF documents are parsed server-side, the extracted
           text is submitted to an LLM for structured analysis, and the resulting
           \texttt{aiVerificationSummary} is surfaced to the administrator. The
           final approval decision is made by a human, preserving oversight in
           this safety-critical process \citep{Dellermann2019}.}
  \label{fig:cred-verification}
\end{figure}

% -----------------------------------------------------------
\subsection{Rate Limiting and API Security}
\label{subsec:rate-limiting}
% -----------------------------------------------------------

In distributed web systems, API endpoints are subject to abuse through
denial-of-service (DoS) attacks, credential stuffing, and excessive automated
requests. Rate limiting is a protective mechanism that restricts the number of
requests a client (typically identified by IP address) can make to an endpoint
within a defined time window.

PsychPlatform applies three distinct rate limiters, as summarised in
Table~\ref{tab:rate-limiters}.

\begin{table}[H]
  \centering
  \caption{Rate limiting configuration per endpoint category in PsychPlatform.}
  \label{tab:rate-limiters}
  \renewcommand{\arraystretch}{1.35}
  \begin{tabularx}{0.85\linewidth}{@{} l X r r @{}}
    \toprule
    \rowcolor{medgray}
    \textbf{Endpoint Category} & \textbf{Purpose} & \textbf{Limit} & \textbf{Window} \\
    \midrule
    Authentication             & Prevent brute-force login attacks & 10 req  & 15 min \\
    \rowcolor{lightgray}
    General API                & Protect backend resources          & 100 req & 1 hour \\
    Chatbot                    & Prevent LLM inference abuse        & 50 msg  & 1 hour \\
    \bottomrule
  \end{tabularx}
\end{table}

% ─────────────────────────────────────────────────────────────
\section{Study of Existing Solutions}
\label{sec:existing}
% ─────────────────────────────────────────────────────────────

This section surveys the landscape of existing platforms and tools that occupy
similar problem spaces to PsychPlatform. The solutions are grouped into five
categories based on primary function and design philosophy.

\begin{figure}[H]
  \centering
  % INSERT FIGURE: Bubble/category map of the existing solutions landscape.
  % Draw five labelled category zones (oval or rounded boxes):
  %   1. Telepsychology Platforms  → BetterHelp, Talkspace, Doctolib
  %   2. AI Mental Health Chatbots → Woebot, Wysa, Replika
  %   3. Video Consultation Tools  → Zoom for Healthcare, Doxy.me, Whereby
  %   4. Practice Management/EHR  → SimplePractice, Theranest
  %   5. AI-Integrated Platforms   → Nabla Copilot, Kry/LIVI
  % Place PsychPlatform at the centre, overlapping all five zones, to
  % visually convey its integrative position.
  \rule{0.85\textwidth}{0.4pt}\\[6pt]
  {\small\itshape [Figure placeholder — Landscape Map of Existing Mental Health Platform Categories]}
  \rule{0.85\textwidth}{0.4pt}
  \caption{Landscape map of the five categories of existing solutions reviewed
           in this chapter. PsychPlatform (centre) is positioned to span all
           five categories, reflecting its integrative design intent: combining
           the discovery features of telepsychology platforms, the AI intelligence
           of chatbot applications, real-time communication, back-office management,
           and AI-assisted clinical workflows.}
  \label{fig:landscape-map}
\end{figure}

% -----------------------------------------------------------
\subsection{Telepsychology and Online Therapy Platforms}
\label{subsec:telepsych}
% -----------------------------------------------------------

This category encompasses commercial platforms whose primary purpose is to
connect patients with licensed mental health professionals over the internet.

\subsubsection{BetterHelp}

BetterHelp is one of the largest online therapy platforms globally, operating in
the United States and internationally. The platform matches patients with
therapists based on a self-reported intake questionnaire covering mental health
history, preferences, and severity of concerns. Communication takes place
through asynchronous text messaging, live chat, video sessions, and phone calls.
Patients can switch therapists upon request, and the platform handles all billing
and scheduling logistics.

\begin{tcolorbox}[colback=lightgray, colframe=accentblue, title=\textbf{BetterHelp — Assessment}, fonttitle=\bfseries, sharp corners=south]
  \textbf{Strengths:} BetterHelp excels in access and scale. It provides therapy
  to patients who would otherwise face geographic, financial, or stigma-related
  barriers. The asynchronous messaging model allows patients to communicate with
  their therapist at any time, which is valuable for patients with non-standard
  schedules or acute moments of distress. \\[4pt]
  \textbf{Limitations:} The platform has faced significant criticism regarding
  practitioner oversight and quality assurance. Critics have noted that the high
  volume of patients per therapist can compromise the depth of therapeutic
  relationships. Furthermore, BetterHelp does not employ any AI integration in
  the therapeutic process; the matching algorithm is opaque and non-adaptive,
  relying entirely on the initial intake questionnaire without learning from
  subsequent interaction patterns.
\end{tcolorbox}

\subsubsection{Talkspace}

Talkspace is a structured teletherapy platform offering text, audio, and video
therapy as well as psychiatry (medication management) services. It differentiates
itself through the provision of psychiatry services and through partnerships with
insurers and employers.

\begin{tcolorbox}[colback=lightgray, colframe=accentblue, title=\textbf{Talkspace — Assessment}, fonttitle=\bfseries, sharp corners=south]
  \textbf{Strengths:} Talkspace's integration with employer health benefits
  programmes and insurance providers represents a significant advancement in
  accessibility. The availability of psychiatry services within a single platform
  allows for more holistic mental health support. \\[4pt]
  \textbf{Limitations:} Like BetterHelp, Talkspace does not employ AI within the
  therapeutic workflow. Session scheduling, therapist matching, and outcome
  tracking are managed manually or through basic algorithmic filtering. The
  platform lacks mechanisms for automated pre-session patient preparation,
  emotional state analysis, or clinical summarisation.
\end{tcolorbox}

\subsubsection{Doctolib (Mental Health Module)}

Doctolib is a dominant European health appointment management platform, with
strong market presence in France, Germany, and Italy. The platform has expanded
to include mental health practitioners, allowing patients to search by specialty,
location, and availability, and book appointments through an integrated calendar
interface.

\begin{tcolorbox}[colback=lightgray, colframe=accentblue, title=\textbf{Doctolib — Assessment}, fonttitle=\bfseries, sharp corners=south]
  \textbf{Strengths:} Doctolib's seamlessly integrated scheduling infrastructure
  and its broad network of certified practitioners are notable. The platform is
  compliant with European health data regulations (GDPR, HDS certification in
  France) and is widely trusted by both patients and clinicians. \\[4pt]
  \textbf{Limitations:} Doctolib is fundamentally a scheduling and administrative
  platform rather than a therapeutic one. It provides no AI-assisted features, no
  inter-session communication, and no mechanisms for patient preparation or
  emotional tracking. It functions as a digital appointment book rather than an
  intelligent healthcare companion.
\end{tcolorbox}

% -----------------------------------------------------------
\subsection{AI-Powered Mental Health Chatbot Applications}
\label{subsec:chatbots}
% -----------------------------------------------------------

This category covers standalone AI chatbot applications designed to provide
psychological support, mood tracking, or Cognitive Behavioural Therapy
(CBT)-based interventions.

\subsubsection{Woebot}

Woebot is a fully automated, AI-powered chatbot application developed by clinical
researchers at Stanford University. It delivers CBT-based therapeutic content
through structured conversational exercises guided by clinical principles and
evidence-based techniques.

\begin{tcolorbox}[colback=lightgray, colframe=accentblue, title=\textbf{Woebot — Assessment}, fonttitle=\bfseries, sharp corners=south]
  \textbf{Strengths:} Woebot is rigorously designed around clinical evidence and
  has been validated in randomised controlled trials demonstrating reductions in
  self-reported anxiety and depression \cite{Fitzpatrick2017}. Its fully
  autonomous design makes it highly scalable and immediately accessible. \\[4pt]
  \textbf{Limitations:} Woebot is a standalone application with no pathway to
  human therapist consultation. The conversational pathways are largely scripted
  according to CBT protocols, with limited ability to deviate from or dynamically
  respond to unexpected patient disclosures.
\end{tcolorbox}

\subsubsection{Wysa}

Wysa is an AI-powered mental health application that combines rule-based CBT and
mindfulness techniques with periodic human coach support. Users interact with an
anonymous AI penguin character and are offered structured coping exercises in
response to expressed emotional states.

\begin{tcolorbox}[colback=lightgray, colframe=accentblue, title=\textbf{Wysa — Assessment}, fonttitle=\bfseries, sharp corners=south]
  \textbf{Strengths:} Wysa's combination of structured AI-driven exercises with
  optional human support represents a thoughtful hybrid model. The anonymised
  character design and conversational tone are intended to reduce stigma and lower
  the barrier to emotional disclosure. \\[4pt]
  \textbf{Limitations:} The human coaching component is limited and text-based.
  The AI component relies on rule-based dialogue management rather than generative
  LLMs, limiting its capacity to handle the full complexity and unpredictability
  of human emotional expression.
\end{tcolorbox}

\subsubsection{Replika}

Replika is a general companionship AI application that employs large language
models to simulate an empathetic conversational companion. While not explicitly
a mental health application, many users turn to Replika for emotional support
and companionship.

\begin{tcolorbox}[colback=lightgray, colframe=accentblue, title=\textbf{Replika — Assessment}, fonttitle=\bfseries, sharp corners=south]
  \textbf{Strengths:} Replika demonstrates the commercial viability and user
  demand for LLM-powered empathetic conversational agents. Its ability to sustain
  long-term, personalised, and emotionally resonant conversations represents a
  significant technical achievement. \\[4pt]
  \textbf{Limitations:} Replika lacks any integration with clinical professionals
  and makes no attempt to route users with serious mental health concerns toward
  professional help. It provides no structured assessment, emotional analysis, or
  clinical summarisation, and its purely companionship-oriented design makes it
  inappropriate for a formal healthcare context.
\end{tcolorbox}

% -----------------------------------------------------------
\subsection{General-Purpose Video Consultation Platforms}
\label{subsec:video}
% -----------------------------------------------------------

This category includes platforms originally designed for video communication
that have been adapted for or adopted in telemedicine and mental health contexts.

\subsubsection{Zoom for Healthcare}

Zoom for Healthcare is a HIPAA-compliant variant of the Zoom video conferencing
platform, used extensively by mental health practitioners in the United States
for conducting teletherapy sessions. It provides end-to-end encrypted video and
audio communication and is compatible with a wide range of devices.

\begin{tcolorbox}[colback=lightgray, colframe=accentblue, title=\textbf{Zoom for Healthcare — Assessment}, fonttitle=\bfseries, sharp corners=south]
  \textbf{Strengths:} Zoom's ubiquity, ease of use, and technical reliability
  make it a practical choice for practitioners who require a simple video channel
  with no additional infrastructure. \\[4pt]
  \textbf{Limitations:} Zoom for Healthcare is a communication tool, not a
  therapeutic platform. It provides no patient management features, no AI
  integration, no scheduling automation, no progress tracking, and no
  inter-session tools. Therapists must manage all administrative and clinical
  functions through separate, disconnected systems.
\end{tcolorbox}

\subsubsection{Whereby and Doxy.me}

Whereby and Doxy.me are browser-based video consultation platforms specifically
designed for medical and therapeutic contexts, operating without software
installation and emphasising simplicity of access for both practitioners and
patients.

\begin{tcolorbox}[colback=lightgray, colframe=accentblue, title=\textbf{Whereby / Doxy.me — Assessment}, fonttitle=\bfseries, sharp corners=south]
  \textbf{Strengths:} The zero-installation, browser-based approach
  significantly reduces friction for first-time patients. Doxy.me has been widely
  adopted by independent therapists for its minimal technical requirements and
  HIPAA-compliant architecture. \\[4pt]
  \textbf{Limitations:} These platforms are communication channels without
  clinical intelligence. They offer no patient history tracking, no AI-assisted
  features, no automated documentation, and no mechanism for session preparation
  or outcome measurement.
\end{tcolorbox}

% -----------------------------------------------------------
\subsection{Practice Management and EHR Systems}
\label{subsec:ehr}
% -----------------------------------------------------------

This category covers software solutions designed for the administrative
management of clinical practice rather than for direct patient engagement.

\subsubsection{SimplePractice}

SimplePractice is a practice management and Electronic Health Records (EHR)
system targeted at solo and small-group mental health practitioners. It provides
scheduling, billing, progress note templates, telehealth integration, and a
patient portal.

\begin{tcolorbox}[colback=lightgray, colframe=accentblue, title=\textbf{SimplePractice — Assessment}, fonttitle=\bfseries, sharp corners=south]
  \textbf{Strengths:} SimplePractice offers a comprehensive administrative suite
  that reduces the documentation burden on practitioners. Its telehealth
  integration and secure messaging features create a more unified experience than
  using disconnected tools. \\[4pt]
  \textbf{Limitations:} SimplePractice is designed primarily as a
  practitioner-side administrative tool. It incorporates no AI features; progress
  notes, treatment plans, and assessments are entirely manual. As a
  practitioner-centric system, it is not designed for patient-initiated discovery
  of psychologists.
\end{tcolorbox}

\subsubsection{Theranest}

Theranest is a similar practice management platform targeting mental health and
counselling practices. It offers appointment scheduling, billing, electronic
signatures, and a client portal.

\begin{tcolorbox}[colback=lightgray, colframe=accentblue, title=\textbf{Theranest — Assessment}, fonttitle=\bfseries, sharp corners=south]
  \textbf{Strengths:} Theranest's group practice features, including staff
  management, treatment team collaboration, and role-based access for multiple
  practitioners, are well-implemented and relevant to larger clinical settings.
  \\[4pt]
  \textbf{Limitations:} Like SimplePractice, Theranest is a back-office system
  with a limited patient-facing interface and no AI integration. It is not a
  public-facing platform through which patients can discover and book sessions
  with new practitioners.
\end{tcolorbox}

% -----------------------------------------------------------
\subsection{AI-Integrated Healthcare and Scheduling Platforms}
\label{subsec:ai-health}
% -----------------------------------------------------------

This emerging category represents platforms that have begun to integrate AI
capabilities into appointment management, clinical documentation, or patient
engagement workflows.

\subsubsection{Nabla Copilot}

Nabla Copilot is an AI-powered clinical documentation assistant that listens to
consultation recordings and automatically generates structured clinical notes
using large language models.

\begin{tcolorbox}[colback=lightgray, colframe=accentblue, title=\textbf{Nabla Copilot — Assessment}, fonttitle=\bfseries, sharp corners=south]
  \textbf{Strengths:} Nabla demonstrates the feasibility and value of
  LLM-based automation of clinical documentation, one of the most time-consuming
  aspects of a clinician's workload. \\[4pt]
  \textbf{Limitations:} Nabla operates as a post-session documentation tool for
  practitioners; it is not patient-facing and does not contribute to inter-session
  support, patient preparation, or emotion tracking. It also does not manage
  sessions, scheduling, or patient-practitioner relationships.
\end{tcolorbox}

\subsubsection{Kry / LIVI}

Kry (marketed as LIVI in the United Kingdom) is a general digital health platform
providing on-demand video consultations with general practitioners and, to a
limited extent, mental health specialists. The platform uses AI assistance in
symptom pre-assessment, asking patients a structured series of questions before
the consultation to provide the clinician with a preliminary patient summary.

\begin{tcolorbox}[colback=lightgray, colframe=accentblue, title=\textbf{Kry / LIVI — Assessment}, fonttitle=\bfseries, sharp corners=south]
  \textbf{Strengths:} Kry's pre-assessment functionality represents a meaningful
  integration of AI into the pre-session workflow, reducing the time practitioners
  spend gathering basic information. Its on-demand model also reduces waiting
  times for patients with episodic needs. \\[4pt]
  \textbf{Limitations:} Kry's AI pre-assessment is form-based and structured,
  lacking the conversational depth and emotional nuance of an LLM-powered chatbot.
  The platform is general-purpose and not optimised for ongoing therapeutic
  relationships or longitudinal mental health monitoring.
\end{tcolorbox}

% ─────────────────────────────────────────────────────────────
\section{Comparative Analysis}
\label{sec:comparison}
% ─────────────────────────────────────────────────────────────

This section provides a structured comparative analysis of the solutions surveyed
in Section~\ref{sec:existing} with respect to a set of functional and technical
dimensions directly relevant to PsychPlatform.

\subsection{Comparison Dimensions}

The following ten dimensions are used for comparison:

\begin{enumerate}[label=\textbf{\arabic*.}, leftmargin=2em, itemsep=3pt]
  \item \textbf{Practitioner Discovery} — Can patients search for and select
        psychologists by specialisation, location, language, and rating?
  \item \textbf{AI Conversational Agent} — Does the platform include an
        LLM-powered conversational chatbot for inter-session patient support?
  \item \textbf{Clinical Summarisation} — Are AI-produced summaries of patient
        interactions made available to the practitioner?
  \item \textbf{Emotional State Tracking} — Does the platform quantitatively
        track patient emotional indicators over time?
  \item \textbf{Credential Verification} — Does credential verification use AI?
  \item \textbf{Real-Time Messaging} — Does the platform support real-time
        text and/or voice messaging?
  \item \textbf{Session Code Security} — Is a secure code-based mechanism used
        for session activation?
  \item \textbf{Multilingual Support} — Does the platform natively support
        multiple languages in both the interface and conversational AI?
  \item \textbf{Geospatial Search} — Is location-based practitioner discovery
        supported?
  \item \textbf{Role-Based Administration} — Is a dedicated administrator role
        provided for oversight and management?
\end{enumerate}

\subsection{Comparison Table}

% Symbols
\newcommand{\yes}{\textcolor{successgreen}{\faCheckCircle}}
\newcommand{\no}{\textcolor{warnorange}{\faTimesCircle}}
\newcommand{\partial}{\textcolor{accentblue}{\faAdjust}}

\begin{table}[H]
  \centering
  \caption{Comparative overview of existing solutions and PsychPlatform across
           key functional dimensions.}
  \label{tab:comparison}
  \small
  \renewcommand{\arraystretch}{1.4}
  \setlength{\tabcolsep}{4pt}
  \begin{tabularx}{\linewidth}{@{} >{\RaggedRight\bfseries}p{2.8cm}
                                   *{10}{>{\centering\arraybackslash}X} @{}}
    \toprule
    \rowcolor{primaryblue}
    \textcolor{white}{\textbf{Solution}}
      & \textcolor{white}{\makecell{\textbf{Disc-}\\\textbf{overy}}}
      & \textcolor{white}{\makecell{\textbf{AI}\\\textbf{Chat}}}
      & \textcolor{white}{\makecell{\textbf{Clin.}\\\textbf{Sum.}}}
      & \textcolor{white}{\makecell{\textbf{Emot.}\\\textbf{Track}}}
      & \textcolor{white}{\makecell{\textbf{Cred.}\\\textbf{Verif.}}}
      & \textcolor{white}{\makecell{\textbf{RT}\\\textbf{Msg}}}
      & \textcolor{white}{\makecell{\textbf{Sess.}\\\textbf{Code}}}
      & \textcolor{white}{\makecell{\textbf{Multi-}\\\textbf{lingual}}}
      & \textcolor{white}{\makecell{\textbf{Geo-}\\\textbf{search}}}
      & \textcolor{white}{\makecell{\textbf{Admin}\\\textbf{Panel}}} \\
    \midrule
    \rowcolor{accentblue!15}
    PsychPlatform (proposed)
      & \yes & \yes & \yes & \yes & \yes & \yes & \yes & \yes & \yes & \yes \\
    BetterHelp
      & \partial & \no & \no & \no & \partial & \yes & \no & \partial & \no & \partial \\
    \rowcolor{lightgray}
    Talkspace
      & \partial & \no & \no & \no & \partial & \yes & \no & \partial & \no & \no \\
    Doctolib
      & \yes & \no & \no & \no & \partial & \no & \no & \partial & \yes & \no \\
    \rowcolor{lightgray}
    Woebot
      & \no & \partial & \no & \partial & N/A & \no & \no & \partial & \no & \no \\
    Wysa
      & \no & \partial & \no & \partial & N/A & \partial & \no & \no & \no & \no \\
    \rowcolor{lightgray}
    Replika
      & \no & \yes & \no & \no & N/A & \no & \no & \partial & \no & \no \\
    Zoom for Healthcare
      & \no & \no & \no & \no & \no & \yes & \no & \yes & \no & \no \\
    \rowcolor{lightgray}
    SimplePractice
      & \no & \no & \no & \no & \no & \partial & \no & \no & \no & \partial \\
    Nabla Copilot
      & \no & \no & \yes & \no & \no & \no & \no & \partial & \no & \no \\
    \rowcolor{lightgray}
    Kry / LIVI
      & \yes & \partial & \partial & \no & \partial & \yes & \no & \yes & \no & \no \\
    \bottomrule
    \multicolumn{11}{@{}l@{}}{\small
      \yes~Full support \quad \partial~Partial / limited \quad \no~Not present \quad N/A~Not applicable}
  \end{tabularx}
\end{table}

\subsection{Analytical Observations}

Several important observations emerge from Table~\ref{tab:comparison}.

\textbf{First}, no existing solution combines all ten dimensions within a single
platform. Current solutions are highly specialised: AI chatbot applications
(Woebot, Wysa) lack practitioner integration; practitioner discovery platforms
(Doctolib, BetterHelp) lack AI intelligence; and clinical documentation tools
(Nabla) lack patient-facing engagement features. PsychPlatform is the only
system in this analysis that integrates all ten dimensions cohesively.

\textbf{Second}, among the platforms that do employ AI, none use it across as
many distinct functional workflows as PsychPlatform. Specifically, PsychPlatform
applies LLM intelligence to:
\begin{enumerate*}[label=(\alph*)]
  \item the conversational agent,
  \item inter-session clinical summarisation,
  \item credential verification,
  \item follow-up question generation for psychologists, and
  \item a platform navigation assistant.
\end{enumerate*}

\textbf{Third}, emotional state tracking — arguably the most clinically valuable
form of AI integration for a mental health platform — is absent from all surveyed
platforms except for the standalone chatbot applications (Woebot, Wysa), and even
those implementations are limited to simple mood ratings rather than the
multi-dimensional, conversation-derived indicators implemented in PsychPlatform.

% ─────────────────────────────────────────────────────────────
\section{Limitations of Current Systems}
\label{sec:limitations}
% ─────────────────────────────────────────────────────────────

Building upon the comparative analysis presented in Section~\ref{sec:comparison},
this section provides a more analytical critique of the structural and systemic
limitations that characterise the current generation of mental health technology
solutions.

\subsection{The Fragmentation Problem}

Perhaps the most pervasive limitation in the current landscape is the
fragmentation of tools across the patient's care journey. A patient who requires
psychological support may currently need to use:
\begin{enumerate}[label=(\arabic*), itemsep=2pt]
  \item a search engine or directory to find a psychologist;
  \item a scheduling platform such as Doctolib to book an appointment;
  \item a separate video conferencing tool for the session itself;
  \item a mental health chatbot for inter-session emotional support; and
  \item a patient portal or email for post-session communication.
\end{enumerate}
This fragmentation imposes significant cognitive and organisational overhead on
patients who are, by definition, in a vulnerable state. It also creates data
silos that prevent any single entity from forming a complete longitudinal picture
of the patient's wellbeing trajectory.

\begin{figure}[H]
  \centering
  % INSERT FIGURE: Fragmentation diagram ("status quo" view).
  % Draw five disconnected boxes representing separate tools the patient must
  % juggle today:
  %   [Directory / Search] → [Scheduling App] → [Video Call Tool]
  %                                  ↓
  %              [Inter-session Chatbot]   [Patient Portal / Email]
  % Connect them with dashed arrows to show the manual hand-offs the patient
  % must perform. Add a DATA SILO annotation under each box to illustrate the
  % absence of a unified data layer. Use a frustrated patient icon at the
  % start of the journey.
  \rule{0.85\textwidth}{0.4pt}\\[6pt]
  {\small\itshape [Figure placeholder — Current Tool Fragmentation in the Patient's Care Journey]}
  \rule{0.85\textwidth}{0.4pt}
  \caption{Illustration of the tool fragmentation currently experienced by a
           patient seeking psychological support. Five disconnected systems are
           involved across the care journey — discovery, scheduling, consultation,
           inter-session support, and post-session communication — each
           maintaining its own isolated data silo and imposing manual
           hand-off overhead on a vulnerable user.}
  \label{fig:fragmentation}
\end{figure}

\subsection{Absence of Pre-Session Patient Preparation}

The majority of platforms surveyed do not address the pre-session period as an
opportunity for clinical value. Patients frequently arrive at therapy sessions
without having reflected systematically on their emotional state, resulting in
sessions that begin slowly or fail to address the most pressing concerns. Neither
BetterHelp, Talkspace, nor Doctolib provide tools that help patients organise
and articulate their thoughts before meeting with their practitioner. The
standalone chatbot applications (Woebot, Wysa) do not bridge into formal
therapeutic relationships and thus do not serve a preparatory function in that
context.

\subsection{Lack of AI-Mediated Clinical Intelligence for Practitioners}

The information available to a psychologist when they begin a session is
typically limited to whatever the patient has chosen to self-report in
registration forms or initial assessments. Platforms do not, in general, provide
the practitioner with AI-generated insight into the patient's recent emotional
state, key preoccupations, or urgency level. This represents a missed opportunity
to improve the efficiency and relevance of clinical consultations. Nabla Copilot
addresses the post-session documentation burden, but no platform reviewed
provides the psychologist with pre-session AI intelligence derived from
inter-session patient interactions.

\subsection{Manual and Opaque Credential Verification}

The verification of practitioner credentials is a critical patient safety
mechanism, yet all surveyed platforms that offer it rely on entirely manual
review processes. This is time-consuming for administrators, creates delays in
onboarding qualified practitioners, and provides no structured assistance to the
reviewing party in assessing the completeness or plausibility of submitted
documentation. The absence of AI-assisted preliminary screening means that the
quality of the review is entirely dependent on the individual reviewer's
expertise and diligence.

\subsection{Limited Multilingual and Cultural Accessibility}

Mental health platforms are predominantly designed for English-speaking markets,
with localisation treated as an afterthought rather than a first-class design
concern. This is particularly problematic in linguistically diverse societies,
where patients who are not fluent in the platform's primary language may be
unable to effectively communicate their emotional state, navigate platform
features, or derive benefit from AI-generated content. Very few platforms support
Arabic, despite the language being spoken by over 400 million native speakers
worldwide. This gap represents both a social equity concern and a significant
commercial opportunity.

\subsection{Absence of Secure, Code-Based Session Authentication}

Online therapy sessions transmit sensitive mental health information that must be
protected from unauthorised access. Most platforms authenticate session access
through standard login credentials alone, which is vulnerable to session
hijacking, credential stealing, and account enumeration attacks. None of the
surveyed platforms employ an out-of-band, one-time session activation code
mechanism to provide an additional layer of assurance that the correct patient is
participating in the session.

\subsection{Static Practitioner Profiles and Insufficient Discovery Mechanisms}

Practitioner discovery on existing platforms is generally limited to keyword
search by name or specialty. Rating and review systems, where they exist, are
either absent or based on superficial star ratings without dimensionality.
Geospatial proximity, which is relevant when patients consider in-person
sessions, is rarely integrated into the search interface. Furthermore,
practitioner profiles are typically static text documents with no dynamic
integration of availability, specialisation filtering, or community-provided
ratings based on multi-dimensional assessment questionnaires.

% ─────────────────────────────────────────────────────────────
\section{Opportunities and Motivation}
\label{sec:opportunities}
% ─────────────────────────────────────────────────────────────

The limitations identified in Section~\ref{sec:limitations} define a clear
opportunity space for an integrated, AI-augmented mental health consultation
platform. This section articulates each opportunity and explains how PsychPlatform
is designed to address it.

\subsection{Towards an Integrated Care Journey}

PsychPlatform responds to the fragmentation problem by integrating the complete
care journey within a single, cohesive platform. A patient can discover a
psychologist, book a session, prepare with an AI chatbot, attend the session
with real-time messaging tools, receive post-session notifications, and rate the
session, all without leaving the platform. This integration eliminates
organisational overhead and ensures that the platform accumulates a rich,
longitudinal dataset about each patient that can be used to continuously improve
the quality of care.

\begin{figure}[H]
  \centering
  % INSERT FIGURE: Unified care journey diagram ("PsychPlatform" view).
  % Draw a single horizontal timeline with six sequential stages, all inside
  % one enclosing "PsychPlatform" boundary box:
  %   [Discover Psychologist] → [Book Session & Pay] → [Chatbot Pre-Session]
  %   → [Live Session (Socket.IO)] → [Post-Session Notification] → [Rate & Review]
  % Below the timeline, draw a shared "Unified MongoDB Data Layer" bar spanning
  % all stages, showing that data flows continuously through one platform.
  % Use icons for each stage (magnifying glass, calendar, chat bubble,
  % video camera, bell, star).
  \rule{0.85\textwidth}{0.4pt}\\[6pt]
  {\small\itshape [Figure placeholder — PsychPlatform Integrated Care Journey Timeline]}
  \rule{0.85\textwidth}{0.4pt}
  \caption{The integrated patient care journey in PsychPlatform. All six stages
           — practitioner discovery, session booking, AI-powered pre-session
           preparation, real-time consultation, post-session notification, and
           multi-dimensional rating — are unified within a single platform backed
           by a shared MongoDB data layer. This eliminates the fragmentation
           and data-silo problem identified in Section~\ref{sec:limitations}.}
  \label{fig:care-journey}
\end{figure}

\subsection{AI-Powered Inter-Session Preparation}

One of PsychPlatform's most distinctive features is its use of an LLM-powered
conversational agent specifically designed to serve as a preparatory bridge
between patient and psychologist. Before a scheduled session, the patient can use
the chatbot to express their feelings, revisit their recent experiences, and
identify the themes they wish to discuss. This conversational activity is not
merely logged; it is analysed by the AI to produce a structured clinical summary
that is made available to the psychologist before the session begins. This
transforms the pre-session period from dead time into a clinically productive
opportunity, and significantly reduces the time required to ``warm up'' a
consultation.

\subsection{AI-Generated Clinical Intelligence for Practitioners}

PsychPlatform addresses the lack of pre-session AI intelligence by providing the
psychologist with a comprehensive, automatically generated summary of the
patient's most recent chatbot interactions. This summary includes the dominant
emotion, sentiment trend, urgency score, key themes, and a narrative summary.
Additionally, the system uses an LLM to generate five tailored follow-up
questions that the psychologist should consider asking during the consultation,
based on the content of the patient's chatbot conversations. This AI-augmented
briefing significantly enhances the practitioner's preparedness and allows the
consultation to be more focused, personalised, and therapeutically effective.

\subsection{AI-Assisted Credential Verification}

The automated credential analysis pipeline implemented in PsychPlatform addresses
the inefficiency of fully manual credential review. By automatically extracting
text from submitted PDFs and submitting that text to an LLM for structured
analysis, the system provides the administrator with a preliminary assessment
summary in seconds rather than hours. While the final approval decision remains
with the human administrator (preserving necessary oversight), the AI
dramatically reduces the cognitive effort required for routine review, enabling
faster onboarding of qualified practitioners and more focused administrative
attention on cases flagged as requiring further scrutiny.

\subsection{Native Multilingual and RTL Support}

PsychPlatform is designed from the ground up as a multilingual platform,
supporting English, French, and Arabic with full right-to-left layout adaptation
for Arabic. This is not a superficial feature: the AI chatbot itself is designed
to detect the patient's language and respond in kind, ensuring that the
highest-value feature of the platform — the therapeutic conversational agent —
is equally accessible to speakers of all three languages. This approach reflects
a commitment to equity of access and positions the platform appropriately for
deployment in North African and Middle Eastern markets, as well as in
French-speaking West African communities.

\subsection{Enhanced Security Through One-Time Session Codes}

PsychPlatform introduces a one-time, six-digit session activation code sent to
the patient's registered email address upon payment confirmation. This code must
be entered before the session becomes active, providing an additional out-of-band
authentication factor that ensures only the registered, paying patient can
initiate the session. This mechanism, absent in all surveyed platforms, mitigates
the risk of session hijacking and provides an auditable payment confirmation
trail.

\subsection{Rich, Quantitative Emotional Longitudinal Tracking}

By maintaining a structured \texttt{EmotionalIndicator} record for each session,
PsychPlatform enables longitudinal analysis of a patient's emotional trajectory
across multiple consultations. Scores for anxiety, sadness, anger, and positivity
are stored per session, creating a time-series dataset that can reveal trends,
regression, or progress in the patient's mental health. This quantitative
tracking provides the psychologist with objective supplementary evidence alongside
their qualitative clinical judgment, supporting more rigorous, outcome-oriented
therapeutic practice.

\begin{figure}[H]
  \centering
  % INSERT FIGURE: Multi-series area/line chart showing emotional trajectory.
  % X-axis: Session number (1 through 8).
  % Y-axis: Score on normalised scale 0-100.
  % Four coloured series: Anxiety (red), Sadness (blue), Anger (orange),
  % Positivity (green). The chart should show a realistic improvement pattern:
  % Anxiety/Sadness/Anger gradually decreasing; Positivity gradually increasing.
  % Add milestone annotations (e.g., "New coping technique introduced" at
  % session 3, "Treatment plan revised" at session 6).
  \rule{0.85\textwidth}{0.4pt}\\[6pt]
  {\small\itshape [Figure placeholder — Longitudinal Emotional Indicator Chart Across Multiple Sessions]}
  \rule{0.85\textwidth}{0.4pt}
  \caption{Illustrative longitudinal \texttt{EmotionalIndicator} chart for a
           hypothetical patient across eight consecutive sessions. Four emotional
           dimensions — anxiety, sadness, anger, and positivity — are tracked
           session-by-session on a normalised scale of 0 to 100. The downward
           trend of negative indicators and the upward trend of positivity reflect
           a therapeutic progression that the assigned psychologist can monitor
           objectively over time.}
  \label{fig:longitudinal-chart}
\end{figure}

\subsection{Multi-Dimensional Practitioner Rating System}

PsychPlatform implements a post-session rating system comprising ten questions
rated on a five-point Likert scale, yielding a maximum score of 50. This
multi-dimensional approach provides a significantly richer and more reliable
measure of patient satisfaction than a simple star rating, enabling meaningful
comparison and quality assurance across the practitioner community. The resulting
average rating and total rating count are displayed on each psychologist's public
profile, supporting informed patient decision-making.

\begin{figure}[H]
  \centering
  % INSERT FIGURE: Feature completeness radar (spider/polar) chart.
  % Ten axes, one per comparison dimension (Practitioner Discovery, AI Chatbot,
  % Clinical Summarisation, Emotional Tracking, Credential Verification,
  % Real-Time Messaging, Session Code, Multilingual, Geospatial, Admin Panel).
  % Plot two overlapping polygons:
  %   - PsychPlatform: filled polygon reaching the outer boundary on all axes.
  %   - Best competing alternative (e.g., Kry/LIVI): a smaller, irregular polygon
  %     covering only some axes.
  % Use a blue fill for PsychPlatform and a grey dashed outline for the comparator.
  \rule{0.85\textwidth}{0.4pt}\\[6pt]
  {\small\itshape [Figure placeholder — Radar Chart: PsychPlatform vs.\ Best Competing Alternative]}
  \rule{0.85\textwidth}{0.4pt}
  \caption{Radar chart comparing PsychPlatform (filled blue polygon) against the
           closest competing solution across the ten functional dimensions defined
           in Section~\ref{sec:comparison}. PsychPlatform achieves full coverage
           on all axes, while existing solutions cover only a subset of dimensions,
           confirming the unique integrative contribution of the proposed platform.}
  \label{fig:radar-chart}
\end{figure}

% ─────────────────────────────────────────────────────────────
\section*{Summary}
\addcontentsline{toc}{section}{Summary}
% ─────────────────────────────────────────────────────────────

\begin{mdframed}[backgroundcolor=lightgray, linecolor=accentblue, linewidth=1.5pt,
                 topline=true, bottomline=true, leftline=false, rightline=false,
                 innertopmargin=10pt, innerbottommargin=10pt]
This chapter has provided a comprehensive review of the scientific domains and
existing solutions that inform and contextualise the development of PsychPlatform.
Section~\ref{sec:domains} established the theoretical foundations of the system,
spanning web engineering, full-stack development, large language models, prompt
engineering, sentiment analysis, geospatial computing, and document-based AI
verification. Section~\ref{sec:existing} presented an in-depth survey of ten
existing platforms across five categories: telepsychology platforms, AI chatbot
applications, general-purpose video consultation tools, practice management
systems, and AI-integrated healthcare solutions. Section~\ref{sec:comparison}
conducted a structured comparative analysis across ten functional dimensions,
demonstrating that PsychPlatform is the only solution in the surveyed landscape
to integrate all dimensions within a single platform. Section~\ref{sec:limitations}
identified seven systemic limitations in existing solutions, including tool
fragmentation, absence of pre-session preparation, lack of AI clinical intelligence
for practitioners, and insufficient multilingual accessibility.
Section~\ref{sec:opportunities} articulated the corresponding opportunities and
explained precisely how each design decision in PsychPlatform is motivated by and
responsive to an identified gap.

The analysis demonstrates that, while individual aspects of PsychPlatform's
functionality are addressed in isolation by existing tools, no single solution
achieves the degree of integration, AI augmentation, clinical intelligence, and
multilingual accessibility that the proposed platform aims to deliver. This
distinction constitutes the primary academic and practical contribution of the
project and motivates the design choices elaborated in the following chapters.
\end{mdframed}

% ─────────────────────────────────────────────────────────────
%  Bibliography entries (add to your .bib file)
% ─────────────────────────────────────────────────────────────
%
% @article{Brown2020,
%   author  = {Brown, T. B. and Mann, B. and Ryder, N. and others},
%   title   = {Language models are few-shot learners},
%   journal = {Advances in Neural Information Processing Systems},
%   volume  = {33},
%   pages   = {1877--1901},
%   year    = {2020}
% }
% @article{Dellermann2019,
%   author  = {Dellermann, D. and Ebel, P. and S{\"o}llner, M. and Leimeister, J. M.},
%   title   = {Hybrid intelligence},
%   journal = {Business \& Information Systems Engineering},
%   volume  = {61},
%   number  = {5},
%   pages   = {637--643},
%   year    = {2019}
% }
% @book{Fedosejev2015,
%   author    = {Fedosejev, A.},
%   title     = {React.js Essentials},
%   publisher = {Packt Publishing},
%   year      = {2015}
% }
% @phdthesis{Fielding2000,
%   author = {Fielding, R. T.},
%   title  = {Architectural Styles and the Design of Network-Based Software Architectures},
%   school = {University of California, Irvine},
%   year   = {2000}
% }
% @article{Fitzpatrick2017,
%   author  = {Fitzpatrick, K. K. and Darcy, A. and Vierhile, M.},
%   title   = {Delivering cognitive behavior therapy to young adults with symptoms of
%              depression and anxiety using a fully automated conversational agent
%              ({Woebot})},
%   journal = {JMIR Mental Health},
%   volume  = {4},
%   number  = {2},
%   pages   = {e19},
%   year    = {2017}
% }
% @book{JurafskyMartin2023,
%   author    = {Jurafsky, D. and Martin, J. H.},
%   title     = {Speech and Language Processing},
%   edition   = {3rd ed. draft},
%   publisher = {Stanford University},
%   year      = {2023}
% }
% @misc{MetaAI2024,
%   author = {{Meta AI}},
%   title  = {Introducing {Meta Llama 3}},
%   year   = {2024},
%   url    = {https://ai.meta.com/blog/meta-llama-3/}
% }
% @book{Mitchell1997,
%   author    = {Mitchell, T. M.},
%   title     = {Machine Learning},
%   publisher = {McGraw-Hill},
%   year      = {1997}
% }
% @incollection{Murugesan2001,
%   author    = {Murugesan, S.},
%   title     = {Web engineering: A new discipline for development of web-based systems},
%   booktitle = {Web Engineering},
%   pages     = {3--13},
%   publisher = {Springer},
%   year      = {2001}
% }
% @misc{OpenAI2023,
%   author = {{OpenAI}},
%   title  = {{GPT-4} Technical Report},
%   year   = {2023},
%   note   = {arXiv:2303.08774}
% }
% @inproceedings{Provos1999,
%   author    = {Provos, N. and Mazi{\`e}res, D.},
%   title     = {A future-adaptable password scheme},
%   booktitle = {Proceedings of the FREENIX Track: USENIX Annual Technical Conference},
%   pages     = {81--91},
%   year      = {1999}
% }
% @book{Russell2020,
%   author    = {Russell, S. and Norvig, P.},
%   title     = {Artificial Intelligence: A Modern Approach},
%   edition   = {4th},
%   publisher = {Pearson},
%   year      = {2020}
% }
% @article{Vaswani2017,
%   author  = {Vaswani, A. and Shazeer, N. and Parmar, N. and others},
%   title   = {Attention is all you need},
%   journal = {Advances in Neural Information Processing Systems},
%   volume  = {30},
%   year    = {2017}
% }
% @misc{W3C2005,
%   author = {{W3C}},
%   title  = {Internationalization ({i18n})},
%   year   = {2005},
%   url    = {https://www.w3.org/International/}
% }
% @misc{White2023,
%   author = {White, J. and Fu, Q. and Hays, S. and others},
%   title  = {A prompt pattern catalog to enhance prompt engineering with {ChatGPT}},
%   year   = {2023},
%   note   = {arXiv preprint arXiv:2302.11382}
% }
% @misc{Anthropic2024,
%   author = {{Anthropic}},
%   title  = {Claude},
%   year   = {2024},
%   url    = {https://www.anthropic.com/}
% }
% @misc{GoogleDeepMind2024,
%   author = {{Google DeepMind}},
%   title  = {Gemini},
%   year   = {2024},
%   url    = {https://deepmind.google/technologies/gemini/}
% }
